\chapter{Experimental Setup}
The experiment is split up into two parts, where the Mach-Zehnder Interferometer is common for them both.
\section{Mach-Zehnder Interferometer}
A helium-neon laser\autocite{thorlabs_henelaser_manual} is used as the light source.
This light is sent through a neutral density filter to decrease its intensity, 
as the sensor\autocite{thorlabs_detector_manual}, oversaturates at the lasers unfiltered intensity.
After moving through the neutral density filter, 
the laserbeam is aligned using two mirrors,
afterwhich it moves through a beamsplitter.
The two split laserbeams then hit a mirror each,
that directs them towards a second beamsplitter,
that unifies the two beams, before they hit the sensor, which is connected to a PicoScope\autocite{PicoScope3203D},
\subsection{Variation with Aircell}
In this setup (see figure \ref{fig:setup_with_aircell}),
an aircell is added to one of the arms of the interferometer. 
This aircell is connected to a bicycle pump that can increase the density of the air inside the aircell (and in turn, the pressure).
\begin{figure*}
    \centering
    \includesvg[width=0.95\textwidth]{../figures/Mach-Zender_IFM_setup_Aircell.svg}
    \caption{Setup of the Mach-Zehnder Interferometer. 
    A HeNe laser beam (red) is attenuated and aligned before entering the first beam splitter. 
    In one of the interferometer arms, the beam through an air cell, where the phase difference between the two arms change due to the higher refractive index ($n$) of the glass and pressurized air.  
    This phase shift is represented as the green colored beam. 
    When the two beams meet again their interference is shown as the dashed overlap of the beams.}
    \label{fig:setup_with_aircell}
\end{figure*}
\subsection{Variation with Waveplate}
This setup is almost the same as before (see figure \ref{fig:setup_with_waveplate}), but in this case the aircell is replaced by a half-waveplate, 
and one of the mirrors is instead mounted on a piezo \autocite{piezo}.
The neutral density filter is also replaced with a polarizer, 
because the polarizer can decrease the intensity, while also allowing us to enforce the polarization of the light.
This polarization angle will later be relevant in calculations of the angle between the fast axis of the waveplate and the polarization of the light.
\begin{figure*}
    \centering
    \includesvg[width=0.95\textwidth]{../figures/Mach-Zender_IFM_setup_HWP.svg}
    \caption{Setup of the Mach-Zehnder Interferometer. 
    A HeNe laser beam (red) is attenuated and aligned before entering the first beam splitter. 
    In one of the interferometer arms, the beam goes thorugh a waveplate, before hitting a mirror mounted on a piezo.
    The phase retarder changes the polarization of the light that moves through it, which is represented as the green colored beam. 
    When the two beams meet again their interference is shown as the dashed overlap of the beams.}
    \label{fig:setup_with_waveplate}
\end{figure*}
\section{Measurements}
\subsection{Variation with Aircell}
First air is pumped into the aircell using the bicycle pump, 
until the gauge pressure exceeds 1 bar. 
The connection between pump and aircell is slightly leaky,
so the pressure will drop slowly. When the pressure is at exactly one bar (measured by eye),
we manually start the measurement. 
The measurement is stopped when the gauge pressure reaches zero (or close to zero, as the pressure changes slower the less air there is in the aircell).
\subsection{Variation with Waveplate}
Firstly we set the polarizer to an angle at which the sensor was not oversaturated. 
This angle ended up being 130 degrees with respect to the vertical.
Then a signal from a signal generator is added to the piezo, so it oscillates back and forth,
changing the path length of one of the light beams, 
and therefore resulting in a signal that oscillates between the maximum and minimum amount of interference.
We do this to ensure that $\Delta \phi$ is always known, as we can always choose the maximum value measured,
and know that at this point $\cos(\Delta \phi)$ must be either $1$ or $-1$ depending on the sign of $\cos(2\theta)$.
If this is not done, we risk that $\Delta \phi$ changes between measurements.
We then rotate the waveplate, 10 degrees at a time, and measure the signal for 5 seconds,
such that the piezo completes several periods. 
We measure the angle of the waveplate using the anglescale on it, 
but it is important to note that the fast axis is offset by 52 degrees, 
meaning that when the angle scale measures 52 degrees, the fast axis is vertical.


